\section{The geometric search model}\label{sec:model}

\subsection{Face-paired triangulations and legal replacements}

A triangulation is a finite list of abstract tetrahedra with local vertices
$0,1,2,3$. Each triangular face is either a boundary face or is paired to
another face by a full vertex permutation. Global vertices and edges may
have multiple local representatives. The existing validator checks the
manifold and matching requirements of the normal-surface interfaces.
This is the usual generalized triangulation setting used in computational
three-manifold topology, rather than a requirement that every global vertex
set specify a unique simplex \cite{BurtonPachner}.

A legal $2$--$3$ replacement consumes two distinct tetrahedra sharing a face.
Their formal union is a triangular bipyramid. It is replaced by the three
tetrahedra around the other diagonal, preserving all six labelled boundary
facets and their attaching permutations. A legal $3$--$2$ replacement consumes
three distinct tetrahedra constituting the entire degree-three star of an
interior edge, with the formal triangular-link consistency checked by the
existing producer. It performs the inverse relative-boundary replacement.

These conventions are substantive. All independence statements below use the
repository's formal-bipyramid legality. If one additionally imposed strict
simplicial-complex conditions forbidding a new edge or face already present
elsewhere, disjoint consumed supports would require a fresh independence
analysis. The source code does not silently switch between these conventions.

\begin{definition}[Incarnation and event]
A tetrahedron incarnation is one specific tetrahedron from its creation until
its destruction. An event is a legal replacement together with the identities
of all consumed incarnations and the local edge or face selected in each.
Surviving incarnations keep their local vertex labels; replacement outputs
receive new identities and the canonical local labels of the producer.
\end{definition}

An incarnation is not a row number. After deleting three rows, a survivor's
row number can change even though its identity, local corners and commitment
remain unchanged. The implementation keeps this distinction explicit.

\subsection{Cocycles and the endpoint surface}

The input marking is a list of four integral heights $h_{t,0},\dots,h_{t,3}$
per tetrahedron. The signed differences on identified edges must agree; they
therefore define an integral cocycle. Adding a constant to one height row
does not change the marking. The implementation normalizes each row by
subtracting its first entry.

For bit-complexity statements, $B$ bounds the bit lengths of both the supplied
height entries and their edge differences. If rows are already normalized,
the edge data alone suffice. Including the supplied entries charges the
input-reading and normalization cost even when a row has an unnecessarily
large additive constant.

If $a,b,c,d$ order the corners so that
$h_a\le h_b\le h_c\le h_d$, the corresponding coherent normal vector has
triangle coordinates $h_b-h_a$ at $a$ and $h_d-h_c$ at $d$, and the single
quadrilateral coordinate $h_c-h_b$ separating $\{a,b\}$ from $\{c,d\}$.
All other local coordinates vanish. Ties cause zero coordinates and cause no
ambiguity in the resulting normal vector.

The five formal bipyramid heights determine the transported cocycle on the
new tetrahedra. This preserves the global cohomology class. It can change the
topology of the coherent normal surface: the incoming transport package gives
a checked example in which a disc becomes a disc and two spheres. Thus raw
Euler characteristic is not used as a terminal unknot criterion.

\begin{definition}[Bounded-upward family]
For a source $(T,h)$ and integer $U\ge0$, let $\Reach_U(T,h)$ be the marked
endpoints of all finite legal $2$--$3$/$3$--$2$ traces containing at most $U$
\emph{total} $2$--$3$ events, with the prescribed cocycle transport at every
event. The empty trace and all intermediate endpoints are included.
\end{definition}

If a trace contains $u$ upward and $d$ downward events, its tetrahedron count is
\begin{equation}\label{eq:tet-account}
 t'=t+u-d.
\end{equation}
In particular, $d\le t+u$ and the trace length is at most $t+2U$. The maximum
tetrahedron count is at most $t+U$. These coarse inequalities also cover
degenerate stopping points without requiring a special minimum-size constant.

The parameter $U$ is different from the maximum excess tetrahedron count and
from the number of upward bursts. Repeated expansion--collapse pairs can have
excess height one and arbitrarily many upward events. Our completeness proof
permutes independent upward and downward events; it preserves total counts,
but does not preserve every prescribed burst pattern.

\begin{center}
\begin{tabular}{@{}ll@{}}
\toprule
Symbol & Meaning\\
\midrule
$n$ & Number of crossings when a knot diagram is the original input\\
$t$ & Number of tetrahedra in the source of one geometric search\\
$U$ & Allowed total number of $2$--$3$ events\\
$F,r$ & Allowed set of original tetrahedra and its size bound\\
$c$ & Maximum number of connected components of the allowed region\\
$B$ & Initial bit-length bound including supplied heights and edge values\\
$Q$ & Number of outer searches in a conditional recognition analysis\\
\bottomrule
\end{tabular}
\end{center}

\subsection{Endpoints, traces and decision predicates}

Two adjacent events with disjoint consumed supports will be shown to commute.
Trace equivalence is generated by such adjacent interchanges, carrying local
labels and the marking through the induced isomorphisms. Endpoint predicates
must respect this transported isomorphism. The existence of an essential-disc
component does; a predicate inspecting the arbitrary order of rows in a
certificate need not.

The code supports an arbitrary callback for experimentation, but
\status{ENDPOINT\_SELECTED} means only that this callback selected a supplied
state. A \status{DISC\_FOUND} result additionally carries an independently
checked component-disc proof. Source geometry is then needed to turn the latter
into \status{UNKNOT} for a given diagram.
